\documentclass[10pt, twoside]{article}
\usepackage{amsfonts,amssymb,latexsym,xspace}
\usepackage{amsmath,enumerate}
\usepackage{pstcol}
\numberwithin{equation}{section}

%\usepackage{showkeys}

\hfuzz 15pt
\newtheorem{thm}{Theorem}[section]
\newtheorem{lem}{Lemma}[section]
\newtheorem{cor}{Corollary}[section]
\newtheorem{sublem}{Sub-Lemma}[section]
\newtheorem{prop}{Proposition}[section]
\newtheorem{defin}{Definition}
\newtheorem{cond}{Condition}
\newtheorem{exmp}{Example}
\newtheorem{rem}{Remark}[section]

\newenvironment{proof}{{\sc Proof\ }}{\hfill $\Box$\vskip.6pt}
\arraycolsep=.5\arraycolsep

\newcommand\Cal[1]{{\mathcal #1}}

\newcommand\B{{\Cal B}}
\newcommand\Co{{\Cal C}}
\newcommand\D{{\Cal D}}
\newcommand\F{{\Cal F}}
\newcommand\Lp{{\Cal L}}
\newcommand\M{{[0,1]}}
\newcommand\T{{\Cal T}}
\newcommand\U{{\Cal U}}
\newcommand\V{{\Cal V}}

\newcommand\Ve{{\mathbb V}}
\newcommand\N{\mathbb N}
\newcommand\R{\mathbb R}
\newcommand\Z{\mathbb Z}
\newcommand\Q{{\mathbb Q}}
\newcommand\E{{\mathbb E}}
\newcommand\Pe{{\mathbb P}}
\newcommand\cp{{\Cal P}}

\newcommand\BV{{\text{BV}}}
\newcommand\vf{\varphi}
\newcommand\nn{\nonumber}
\newcommand\ve{\varepsilon}

%%%%%   Figure-picture related definitions--beginning
\newenvironment{myfigure}{\begin{figure}[ht]\
\centering}{\end{figure}}

\definecolor{lgray}{gray}{0.9}
\definecolor{gray}{gray}{0.8}
\definecolor{dgray}{gray}{0.7}
\setlength{\unitlength}{1mm}
%%%%%   Figure related definitions--end

\pagestyle{myheadings}
\markboth{\centerline{\sc Carlangelo Liverani}\hskip-12pt}
{\centerline{\sc problemi concernenti i limiti }\hskip-12pt}

\begin{document}
\title{\bfseries Alcuni problemi non banali concernenti i limiti}
\author{{\sc Carlangelo Liverani}\\
Dipartimento di Matematica\\
Universit\`{a} di Roma (Tor Vergata)\\
Via della Ricerca Scientifica, 00133 Roma, Italy\\
{\tt liverani@mat.uniroma2.it}}
\date{Rome, Novembre 5, 2001}
\maketitle


\section{\sc \large Quanto spesso vi pagano gli interessi?}

Supponiamo di avere un capitale  di $a>0$ euro e di metterlo in una
banca che ci paga l'uno per cento annuo su base annua. Questo significa che, dopo un
anno, il nostro capitale sar\`a $a+a/100=a(1+1/100)$. A questo punto riceviamo
una pubblicit\`a di un'altra banca che di offre lo stesso tasso di
interesse ma che lo paga su di una base semestrale. Ci\`o significa
che dopo sei mesi paga mezzo punto, dunque dopo sei mesi il nostro
capitale sar\`a $a+a/200$ e dopo altri sei mesi ci vene pagato
l'interesse su questo capitale, dunque dopo un anno avremo
\[
\left(a+\frac a{200}\right)+\frac 1{200}\left(a+\frac a{200}\right)=
a\left(1+\frac 1{200}\right)^2.
\]
Per esempio, se il nostro capitale iniziale \`e di $10000$ euro
($a=10000$),
allora nel primo caso dopo un anno abbiamo $10100$ euro, nel secondo
$10100.25$ euro (arrontondato al centesimo). Insomma, non molto di
pi\`u ma comunque di pi\`u. Stiamo per trasferire i nostri fondi in
tale banca quando veniamo a sapere un'altra banca paga
gli intreressi su base quadrimestrale, cio\`e tre volte all'anno. Un
rapido calcolo ci fa vedere che in questo caso alla fine di un anno il
nostro capitale sarebbe
\[
a\left(1+\frac 1{300}\right)^3.
\]
Dunqe, se $a=10000$ si ottiene$10100.33$ (arrotondato al centesimo),
ancora un poco di pi\`u che nel caso presedente. Saputo ci\`o cadiamo
vittima dell'avidit\`a e ci chiediamo che succederebbe se
ci pagassero gli interessi $n$ volte all'anno. Chiaramente dopo un
anno avremmo
\begin{equation}\label{eq:standard}
a\left(1+\frac 1{100 n}\right)^n.
\end{equation}
Che succede a questa successione, cresce arbitrariamente oppure,
purtroppo, \`e limitata?

Per capirlo vediamo un caso pi\`u semplice, ovvero dallo studio della
successione\footnote{Si tratterebbe di una banca assurdamente generosa, che
paga un interesse del 100\%.}
\[
a_n=\left(1+\frac 1{n}\right)^n.
\]
Ci serve un risultato ausiliario.
\begin{lem}\label{lem:semindo}
Per ogni $n\in\N$ ed ogni $0\leq k\leq n/2$ si ha
\[
\left(1+\frac 1{n}\right)^k\leq 1+\frac{2k}n.
\]
\end{lem}
\begin{proof}
Cominciamo col vedere che succede se $n=1$ e qundi $k=0$.
\[
(1+1)^0= 1\leq 1
\]
dunque in qusto caso il lemma \`e vero. Proviamo con $n=2$, chiaramente le sole
possibilit\`a sono $k\in\{0,1\}$. Se $k=0$ si ha $1\leq 1$, se $k=1$
\[
(1+1/2)^1=3/2\leq 1+2/2=2,
\]
e nuovamente la cosa \`e vera. Ovviamente, poich\`e il numero dei casi
da verificare \`e infinito non possiamo continuare in questo modo, \`e
necessaria una strategia migliore, questa \`e fornita dal
\emph{principio di induzione}.

Consideriamo un $n$ arbitrario allora per $k=0$ abbiamo $1\leq 1$, per
$k=1$
\[
(1+1/n)^1\leq 1+2/n
\]
che ancora \`e vero. Supponiamo ora di avere controllato la validit\`a
della affermazione per tutti i $k$ pi\`u piccoli di un certo
$k_*<n/2$. Consideriamo allora $k_*+1$. Se $k_*+1>n/2$ allora abbiamo
finito, altrimenti
\begin{equation}\label{eq:step}
\begin{split}
(1+1/n)^{k_*+1}=&(1+1/n)^{k_*}(1+1/n)\leq
\left(1+\frac{2k_*}n\right)(1+1/n)\\
=&1+\frac{2k_*+1}{n}+\frac{2k_*}{n^2}
\leq1+\frac{2(k_*+1)}{n}.
\end{split}
\end{equation}
Da cui segue che l'affermazione del Lemma deve essere valida anche per
$k_*+1$. Ma allora, usando questo risultato, e poich\`e abbiamo gi\`a
controllato il caso $k_*=1$, ne segue che l'affermazioen deve essere
valida per $k_*+1=2$. Ma ora nulla ci impedisce di usare
(\ref{eq:step}) con $k_*=2$ ottendo perci\`o la validit\`a del Lemma
per $k=3$. Chiaramente continuando in questo modo dopo un numero
finito di passi si sar\`a controllato la validit\`a del Lemma per
tutti i $k\leq n/2$. Il Lemma segue allora dalla arbitrariet\`a di $n$.
\end{proof}
Usando il risultato precedente siamo in grado di vedere che la
successione $a_n$  \`e limitata. Infatti se $n$ \`e pari ($n=2m$)
possiamo scrivere
\[
\left(1+\frac 1{n}\right)^n=\left[\left(1+\frac
1{n}\right)^m\right]^2\leq
\left(1+\frac{2m}{n}\right)^2=4.
\]
Se invece $n$ \`e dispari ($n=2m+1$) allora
\[
\left(1+\frac 1{n}\right)^n=\left[\left(1+\frac
1{n}\right)^m\right]^2\left(1+\frac 1{n}\right)\leq 4\left(1+\frac
1{n}\right)\leq 8.
\]

A questo punto per vedere che esiste il limite basterebbe l'andamento
che abbiamo visto all'inizio di 
questa discussione, ossia che la successione \`e crescente. Per vedere
ci\`o \`e necessario un altro risultato ausiliario.\footnote{Il
seguente risultato \`e noto come \emph{disuguaglianza di Bernoulli}.}
\begin{lem}\label{lem:bernu}
Per ogni $b>-1$ ed ogni $n\in\N$ si ha
\[
(1+b)^n\geq 1+nb.
\]
\end{lem}
\begin{proof}
Se $n=1$ si ha una uguaglianza, dunque il Lemma \`e vero per $n=1$. A
questo punto usiamo di nuovo il principio di
induzione.\footnote{Questa volta si tratta di una cosa pi\`u seria
visto che lo usiamo non sull'insieme finito $k\leq n/2$, bens\`\i\
sull'insieme infinito $\N$.} Supponaimo il Lemma vero per un qualche
$n\in\N$ allora
\[
(1+b)^{n+1}=(1+b)^n(1+b)\geq (1+nb)(1+b)=1+(n+1)b+nb^2\geq 1+(n+1)b.
\]
La prima disuguaglianza segue dall'ipotesi che $1+b>0$ e dall'ipotesi
induttiva $(1+b)^n>1+nb$.
Dunque, poich\`e \`e vero per uno, sar\`a vero per due, per tre e
cos\`\i\ via.
\end{proof}
A questo punto per verificare che la successione \`e crescente
dobbiamo dimostrare che per ogni $n$ vale $a_{n+1}\geq a_n$. 
\[
\begin{split}
\frac {a_{n+1}}{a_n}=&\left(1+\frac 1{n+1}\right)\left[\frac {1+\frac
1{n+1}}{1+\frac 1n}\right]^n=\left(1+\frac 1{n+1}\right)\left(1-\frac
1{(n+1)^2}\right)^n \\
\geq& \left(1+\frac 1{n+1}\right)\left(1-\frac n{(n+1)^2}\right)
\geq 1
\end{split}
\]
dove abbiamo usato Lemma \ref{lem:bernu} con $b=-n(n+1)^{-2}$.

Dunque la successione (\ref{eq:standard}) \`e monotona crescente e
limitata e quindi ammette limite, chiamiamo $e$ tale limite. Notare
che, sebbene abbiamo stabilito la sua esistenza, non sappiamo un gran
ch\`e sul numero $e$.
A questo punto possiamo tornare alla nostra successione originaria
\[
\lim_{n\to\infty}\left(1+\frac{1}{100n}\right)^{n}=
\lim_{n\to\infty}\left[\left(1+\frac{1}{100n}\right)^{100n}\right]^{\frac
1{100}}=e^{\frac 1{100}}.
\]

\section{\sc \large L'area del cerchio}

In questa sezione cercheremo di calcolare l'area del cerchio seguendo
una vecchia idea di Archimede (III secolo avanti Cristo).

Si consideri un cerchio di raggio $r$, chiaramente la sua area \`e
pi\`u grande di quella di un qualunque poligono inscritto e pi\`u
piccola di quella di un qualunque poligono circoscritto. Ispirandoci ad
Archimede considereremo solo poligoni regolari con $n$ lati. Si
veda in figura \ref{fig:pol} uno spicchio di poligono inscritto e
circoscritto di angolo $2\alpha$, cio\`e con $\pi/\alpha$ lati
(ovviamente supponendo $\pi/\alpha$ intero).
\begin{myfigure}
\put(-30,0){\line(1,0){60}}
\put(0,-30){\line(0,1){60}}
\pscircle(0,0){2}
\put(0,0){\line(2,1){20}}
\put(0,0){\line(2,-1){20}}
\put(20,-10){\line(0,1){20}}
\put(17.7,-8.8){\line(0,1){17.5}}
\put(6,1){$\alpha$}
\caption{Uno spicchio di poligono inscritto e di poligono
circoscritto}\label{fig:pol} 
\end{myfigure}
Vediamo l'area del poligono inscritto. Se il poligono ha $n$ lati
allora $\alpha=2\pi/n$. L'area del mezzo spicchio in figura \`e
$r^2/2\cdot \sin\alpha\cos\alpha=r^2/4\cdot \sin2\alpha$. Chiamando
$A_{i,n}$ l'area del poligono inscritto di $n$ lati abbiamo perci\`o
\[
A_{i,n}=r^2\frac n 2\sin\left(\frac {2\pi}n\right).
\]
Analogamente l'area del poligono circoscritto $A_{c,n}$ \`e data da
\[
A_{c,n}=r^2 n \tan\left(\frac\pi n\right).
\]
Chiaramente, per ogni $n\in\N$
\begin{equation}\label{eq:conf}
A_{i,n}\leq A_c\leq A_{c,n}
\end{equation}
dove $A_c$ \`e l'area della circonferenza.

Studiamo cosa succede alle espressioni di cui sopra quando $n$ tende
all'infinito. 
Dalla figura \ref{fig:pol} \`e evidente che per ogni angolo $\alpha$
si ha
$\sin\alpha\leq \alpha\leq \tan\alpha$, poich\`e $r\alpha$ \`e la
lunghezza dell'arco. Da questo segue $\alpha\cos\alpha\leq
\sin\alpha\leq \alpha$ ovvero
\begin{equation}\label{eq:basic}
\cos\alpha\leq \frac{\sin\alpha}{\alpha}\leq 1.
\end{equation}
Ponendo $\alpha=\frac{2\pi}n$ in (\ref{eq:basic}) otteniamo
\[
\cos \frac {2\pi}n\leq\frac{\sin\left(\frac
{2\pi}n\right)}{\frac {2\pi}n} \leq 1
\]
ovvero
\[
\pi r^2\cos \frac {2\pi}n\leq A_{i,n}\leq \pi r^2.
\]
Visto che chiaramente $\lim\limits_{n\to\infty}\cos \frac {2\pi}n=1$,
ne segue, per il confronto, che $\lim\limits_{n\to\infty}A_{i,n}=\pi
r^2$.

Analogamente, da (\ref{eq:basic}),
\[
1\leq \frac{\tan\alpha}{\alpha}\leq \frac 1{\cos\alpha}
\]
e, allo stesso modo che sopra, $\lim\limits_{n\to\infty}A_{c,n}=\pi
r^2$.

Finalmente, prendendo il limite per $n$ che va all'infinito nella
\ref{eq:conf} si ottiene la ben nota formula:
\[
A_c=\pi r^2.
\]

Notare che tutto si basa sul limite notevole
\[
\lim_{x\to 0}\frac {\sin x}{x}=1,
\]
la cui validit\`a \`e sopra dimostrata.

\section{\sc \large Il logaritmo}

L'utilit\`a del logaritmo nasce dalla sua caratteristica di
trasformare prodotti in somme facilitando in tal modo calcoli
complessi.\footnote{Con l'avvento dei calcolatori questo aspetto dei
logaritmi ha perso di importanza ma non per questo l'utilit\`a dei
logaritmi \`e venuta a cessare.}
Vediamo come nasce una tale funzione.

Supponiamo di cercare una funzione $\ell:\R\to\R$ che abbia la magica
propriet\`a 
\begin{equation}\label{eq:magic}
\ell(ab)=\ell(a)+\ell(b),
\end{equation}
esiste una tale funzione e se si come \`e fatta?

Prima di tutto notiamo che per ogni numero $a\in\R$, dovrebbe essere
\[
\ell(0)=\ell(0\cdot a)=\ell(0)+\ell(a) 
\]
ma questo implicherebbe $\ell(a)=0$ sempre, una funzione non molto
interessante! Rassegnamoci dunque al fatto che una tale funzione non
pu\`o essere definita in $0$. Limitamoci a cercare
$\ell:\R^+\backslash\{0\}\to \R$ 
con la propriet\`a (\ref{eq:magic}).
\[
\ell(1)=\ell(1\cdot 1)=2\ell(1)
\]
dunque deve essere $\ell(1)=0$. D'altro canto
\[
\ell(a^p)=p\ell(a)
\]
e
\[
0=\ell(1)=\ell(aa^{-1})=\ell(a)+\ell(a^{-1})
\]
dunque $\ell(a^{-1})=-\ell(a)$. Infine,
\[
\ell(a)=\ell(a^{\frac qq})=q\ell(a^{\frac 1q})
\]
cio\`e $\ell(a^{\frac 1q})=\frac 1q\ell(a)$.
In conclusione abbiamo che, per ogni numero $b\in\Q$ si ha
\begin{equation}\label{eq:mainraz}
\ell(a^b)=b\ell(a).
\end{equation}
Questa propriet\`a \`e ovviamente molto conveniente per calcolare
potenze e radici di numeri. se $\bar a\in\R^+$ \`e tale $\ell(\bar
a)=1$ allora la (\ref{eq:mainraz}) diventa $\ell(a^b)=b$; cio\`e sui
razionali $\ell$ \`e l'inverso della elevazione a potenza. 
Inoltre \`e chiaro che se una funzione
$\ell$ soddisfa (\ref{eq:magic}) allora anche la funzione $\lambda
\ell$, per ogni $\lambda \in\R$ soddisfer\`a  (\ref{eq:magic}).

Questo ci da una buona idea delle propriet\`a di $\ell$ ma la domanda
rimane: esiste una tale funzione sui reali con la magica propriet\`a
(\ref{eq:magic})?

Per rispondere a questa ultima domanda \`e possibile ragionare in
diversi modi. Il mod sotto riportato \`e interessante per il suo
particolare sapore geometrico.

Purtroppo \`e necessario iniziare con una piccola digressione.

Per ogni $a\in\R^+$ si consideri la trasformazione del piano
$T_a(x,y):=(a x, a^{-1}y)$. Se prendiamo il quadrato $Q=\{(x,y)\in
\R^2\;|\; x\in [1,2],y\in[0,1]\}$ allora $T_aQ$ \`e il rettangolo $\{(x,y)\in
\R^2\;|\; x\in [a,2a],y\in[0,a^{-1}]\}$, si veda la figura
\ref{fig:transf} per un esempio della azione di $T_a$ nel caso $a=2$.

\begin{myfigure}
\put(0,33){\ }
\psframe[fillstyle=solid, fillcolor=lgray](-4,0)(-1,3)
\psframe[fillstyle=solid, fillcolor=gray](-1,0)(4,1.5)
\put(-41,-4){$1$}
\put(-11,-4){$2$}
\put(39,-4){$4$}
\put(-30,15){$A$}
\put(20,6){$T_2A$}
\caption{Azione di $T_{2}$.}\label{fig:transf}
\end{myfigure}

La cosa interessante delle trasformazioni $T_a$ \`e che l'immagine di
un rettangolo \`e un rettangolo e che l'area del rettangolo immagine
\`e uguale all'area del rettangolo originario (uno nell'esempio in figura).
Sembra quindi naturale che questa propriet\`a di conservare le aree
valga per qualunque (ragionevole) regione del piano. Assumiamolo e vediamo che
ne segue.\footnote{Il problema chiaramente \`e che non abbiamo
ancora definito l'area di una figura \emph{qualsiasi} e dunque non \`e
ben chiaro come calcolarne algebricamente il valore,
tuttavia il suo significato geometrico--intuitivo dovrebbe essere chiaro.}

Si consideri la funzione $f(x)=\frac 1 x$ per $x>0$, si tratta di una
iperbole. Si definisca la
funzione $\ln:\R^+\to\R$ nel modo seguente: se $z\geq 1$ allora $\ln(z)$
\`e uguale all'area della regione $P(z)=\{(x,y)\in\R^2\;|\; x\in[1,z];
0\leq y\leq 1/x\}$. Per esempio l'area della regione grigio chiaro nella
figura \ref{fig:inv} \`e il valore di $\ln(2)$. Se invece $z<1$ allora
$\ln(z)$ \`e uguale all'area della regione
$P(z)=\{(x,y)\in\R^2\;|\; x\in[z,1]; 0\leq y\leq 1/x\}$ presa col
segno negativo. Chiaramente
$\ln(z)$ \`e una funzione ben definita (almeno dal punto di vista
geometrico). 


\begin{myfigure}
\put(-60,0){\line(1,0){120}}
\put(-60,0){\line(0,1){60}}
\pscurve(-4.5,6)(-4,4.5)(-3,3)(-2,2.25)(-1,1.8)(0,1.5)(1,1.2857)(2,1.125)%
(3,1)(4,.9)(5,.8182)(6,.75)
\put(-30,-5){$1$}
\pscustom[fillstyle=solid,fillcolor=lgray]{%
\psecurve(-4,4.5)(-3,3)(-2,2.25)(-1,1.8)(0,1.5)(1,1.2857)
\psline[linestyle=dashed](0,1.5)(0,0)
\psline[linewidth=.2pt, linestyle=dashed](-3,0)(-3,3)
}
\put(0,-5){$2$}
\put(60,-5){$4$}
\put(-60,-5){$0$}
\pscustom[fillstyle=solid,fillcolor=dgray]{%
\psecurve(-1,1.8)(0,1.5)(1,1.2857)(2,1.125)%
(3,1)(4,.9)(5,.8182)(6,.75)(7,.6923)
\psline[linestyle=dashed](6,.75)(6,0)
\psline[linewidth=.2pt, linestyle=dashed](0,0)(0,1.5)
}
\caption{Grafico di $f(x)=1/x$}\label{fig:inv} 
\end{myfigure}

Vediamo che succede applicando la trasformazione $T_a$ alla regione
$P(z)$. Se $y\geq 0$ allora $a^{-1}y\geq 0$ inoltre se $y\leq 1/x$ allora
$a^{-1}y\leq a^{-1}/x=1/(ax)$. Ci\`o significa che se un punto
appartine all'iperbole anche il punto immagine apparterr\`a
all'iperbole, mentre se un punto si trova sotto l'iperbole allora
anche il punto immagine giace sotto l'iperbole, 
dunque $T_aP(z)=\{(x,y)\in\R^2\;|\;x\in[a,za]; 0\leq y\leq 1/x\}$. Nella
figura \ref{fig:inv} si vede la regione $T_aP(z)$, nel caso speciale in
cui $z=2$ ed $a=2$, obreggiata in grigio scuro.

Da quanto detto segue che $P(a)\cup T_aP(b)=P(ab)$ (nel caso in figura
$P(2)\cup T_2P(2)=P(4)$). Ma questo significa
\[
\ln(ab)=\ln(a)+\ln(b)
\]
cio\`e la funzione $\ln$ ha la propriet\`a (\ref{eq:magic})! Dunque tali
funzioni esistono, ma siamo in grado di calcolarne i valori?


Chiaramente $\ln$ \`e crescente ed illimitata, dunque \`e ragionevole
assumere che esista un numero reale $e$ tale che $\ln(e)=1$, vediamone
comunque una dimostrazione rigorosa basta su di un modo di ragionare
che ci dovrebbe essere ormai familiare.
\begin{lem}\label{eq:e}
Existe un numero $e\in\R$ tale che $\ln(e)=1$.
\end{lem}
\begin{proof}
Prima di tutto vediamo che esiste un munero $a_1$ tale che
$\ln(a_1)>1$. Poich\`e dalla figura \ref{fig:inv} \`e ovvio $\ln 2>0$
ne segue che $\ln 2^n=n\ln 2>1$ a patto che $n$ sia sufficientemente
grande (pi\`u grande di $(\ln 2)^{-1}$ per la precisione). Abbiamo
perci\`o che nell'intervallo $[1,a_1]$ la funzione $\ln$ cresce da
$0$ ad un valore maggiore di $1$. Dividiamo l'intervallo a met\`a (sia
$a_2$ il punto medio) se $\ln a_2=1$ abbiamo trovato il numero che
cercavamo, se $\ln a_2<1$ allora consideriamo l'intevallo $[a_2,a_2]$;
se invece $\ln a_2>1$ allora consideramo l'intervallo $[1,a_2]$. In
entrambi i casi abbiamo un intervallo di lunghezza $(a_1-1)/2$ in cui
$\ln$ vale meno di uno nell'estremo di sinistra e vale pi\`u di uno
sull'estremo di destra (comme sull'intervallo orinario). Ora dividiamo
nuovamente l'intervallo a met\`a, chiamiamo $a_3$ il punto medio, e
ragionando come sopra otteniamo un nuovo intervallo, questa volta di
lunghezza  $(a_1-1)/4$, tale che $\ln$ vale meno di uno nell'estremo
di sinistra e vale pi\`u di uno sull'estremo di destra (a meno che
$a_3$ non fosse il numero cercato).

A questo punto possiamo iterare questa procedura quante volte vogliamo
e, se non troviamo prima il numero voluto, otteniamo una successione
$\{a_n\}$, chiaramente tale successione \`e di Cauchy, sia $e$ il suo
limite. Per costruzione  
\[
|e-a_n|\leq \frac{a_1-1}{2^{n-1}}.
\]
Rimane da investigare che succede alla successione $\{\ln a_n\}$.
Se $a>b>1$ allora $\ln a-\ln b\leq b^{-1}(a-b)$, questo perch\`e $\ln
a-\ln b$ \`e semplicemente l'area sotto l'perbole nell'intervallo
$[b,a]$ mentre $ b^{-1}(a-b)$ \`e l'area di un rettangolo di base
$[b,a]$ e altezza $b^{-1}$ che ovviamente contiene l'area precedente
(si veda la figura \ref{fig:rects0}).
\begin{myfigure}
\put(-30,0){\line(1,0){60}}
\put(-30,0){\line(0,1){60}}
\psframe[fillstyle=solid, fillcolor=white,linewidth=1pt](0,0)(1,3)
\pscurve (-1.5,6)(-1,4.5)(0,3)(1,2.25)(2,1.8)(3,1.5)
\pscustom[fillstyle=solid,fillcolor=gray]{%
\psecurve(-1,4.5)(0,3)(1,2.25)(2,1.8)
\psline[linestyle=dashed](1,2.25)(1,0)
\psline[linewidth=.2pt, linestyle=dashed](0,0)(0,3)
}
\put(0,-5){$b$}
\put(10,-5){$a$}
\multiput(-30,30)(2,0){15}{\line(1,0){1}}
\put(-36,29){$b^{-1}$}
\caption{Come stimare l'area sotto l'iperbole}\label{fig:rects0}
\end{myfigure}
Dunque si ha $|\ln a_n-\ln e|\leq |e-a_m|\leq
(a_1-1)2^{-n+1}$, perci\`o $\ln e=\lim\limits_{n\to\infty} a_n$. Ma 
$\lim\limits_{n\to\infty} a_n= 1$ per costruzione e questo conclude il
Lemma. Infatti, la successione $\{\ln a_n\}$ \`e l'unione di due 
sottosuccessioni: quella degli $a_n$ che sono il limite sinistro di
uno degli intervalli che abbiamo costruito iterativamente e quella
degli $a_n$ che sono un limite destro. Queste due sottosuccessioni
devono avere lo stesso limite della successione $\ln a_n$ ma una deve
avere limite minore od uguale ad uno e l'altra maggiore o uguale a uno, per
il modo con cui sono stati costruiti gli intervalli. Dunque l'unica
possibilit\`a \`e che il limite sia uno.  

\end{proof}
Dunque il numero $e$ esiste, ma di che
numero si tratta?

Consideriamo $\ln(1+1/n)$, chiaramente
\begin{equation}\label{eq:mainineq}
\frac 1{1+\frac 1n}\frac 1n=\frac 1{n+1}\leq ln\left(1+\frac
1n\right)\leq \frac 1n. 
\end{equation}
Infatti guradando alla figura \ref{fig:rects}, si vede che la prima
quantit\`a \`e l'area del rettangolo grigio scuro, la seconda (in
grigio) \`e l'area sotto l'iperbole (che definisce $\ln$)
e la terza \`e l'area del rettango pi\`u alto (in grigio chiaro).
\begin{myfigure}
\put(-30,0){\line(1,0){60}}
\put(-30,0){\line(0,1){60}}
\psframe[fillstyle=solid, fillcolor=lgray,linewidth=1pt](0,0)(1,3)
\pscurve (-1.5,6)(-1,4.5)(0,3)(1,2.25)(2,1.8)(3,1.5)
\pscustom[fillstyle=solid,fillcolor=gray]{%
\psecurve(-1,4.5)(0,3)(1,2.25)(2,1.8)
\psline[linestyle=dashed](1,2.25)(1,0)
\psline[linewidth=.2pt, linestyle=dashed](0,0)(0,3)
}
\psframe[fillstyle=solid, fillcolor=dgray,linewidth=1pt](0,0)(1,2.25)
\put(0,-5){$1$}
\put(8,-5){$1+\frac 1n$}
\caption{Aree rilevanti per l'equazione (\ref{eq:mainineq})}\label{fig:rects}
\end{myfigure}

Dalla (\ref{eq:mainineq}) e dalle propriet\`a di $\ln$ segue che
\begin{equation}\label{eq:este1}
1\leq \ln\left(1+\frac
1n\right)^{n+1}\quad \text{ e che }\quad \ln\left(1+\frac
1n\right)^n\leq 1.
\end{equation}
Poich\`e la funzione $\ln$ \`e crescente
\begin{equation}\label{eq:neper}
\left(1+\frac 1n\right)^n\leq e\leq \left(1+\frac 1n\right)^{n+1}.
\end{equation}
Vediamo dunque che $e$ non era un nome casuale, si tratta infatti
della nostra vecchia conoscenza (si veda la prima sezione di queste note).
D'altro canto le (\ref{eq:este1}) implicano
\[
\left(1+\frac 1n\right)^{n+1}- \left(1+\frac 1n\right)^n\leq \frac en
\]
e da questo segue una dimostrazione alternativa del limite notevole
\[
\lim_{n\to\infty}\left(1+\frac 1n\right)^n=e.
\]
Per di pi\`u la (\ref{eq:mainineq}) ci permette una conoscenza pi\`u
precisa del misterioso numero $e$, infatti prendento, ad esempio,
$n=5$ si ottiene $2.48832\leq e\leq 2.985984$.\footnote{Il lettore che
ha voglia di fare alcuni calcoli numerici pu\`o facilmente ottenere
stime migliori scegliendo un $n$ pi\`u grande.}
\end{document}

